% ============================================================
%  CHAPITRE 1 : CADRE GÉNÉRAL
% ============================================================
\chapter{Cadre général}

\section*{Introduction}
\addcontentsline{toc}{section}{Introduction}
Ce chapitre présente le cadre général du projet ainsi que le contexte de développement de la plateforme FENDRI. Il expose les objectifs du projet, la problématique identifiée ainsi qu'une étude de l'existant permettant de mieux comprendre les besoins et les limites des solutions actuelles.

% -----------------------------------------------------------
\section{Présentation de l'organisme d'accueil}
% -----------------------------------------------------------
% [Copier le texte depuis Canva]

\begin{figure}[H]
  \centering
  \includegraphics[width=4cm]{logo_digilab.png}
  \caption{Logo de DigiLab Solutions}
  \label{fig:logo_digilab}
\end{figure}

% -----------------------------------------------------------
\section{Contexte général du projet}
% -----------------------------------------------------------
% [Copier le texte depuis Canva]

% -----------------------------------------------------------
\section{Description du projet}
% -----------------------------------------------------------
% [Copier le texte depuis Canva]

% -----------------------------------------------------------
\section{Objectifs à atteindre}
% -----------------------------------------------------------

\subsection{Objectif général}
% [Copier le texte depuis Canva]

\subsection{Objectifs spécifiques}
% [Copier le texte depuis Canva]

% -----------------------------------------------------------
\section{Problématique}
% -----------------------------------------------------------
% [Copier le texte depuis Canva]

% -----------------------------------------------------------
\section{Étude de l'existant}
% -----------------------------------------------------------

\subsection{Analyse de l'existant}

% ---- Plateforme 1 : Carthage Crown ----
\subsubsection*{Carthage Crown}
% [Description]

\begin{figure}[H]
  \centering
  \includegraphics[width=\textwidth]{carthage_crown_accueil.png}
  \caption{Interface d'accueil de la plateforme Carthage Crown}
  \label{fig:carthage_crown_accueil}
\end{figure}

\begin{figure}[H]
  \centering
  \includegraphics[width=\textwidth]{carthage_crown_produits.png}
  \caption{Présentation des produits sur la plateforme Carthage Crown}
  \label{fig:carthage_crown_produits}
\end{figure}

% ---- Plateforme 2 : ZETUN ----
\subsubsection*{ZETUN}
% [Description]

\begin{figure}[H]
  \centering
  \includegraphics[width=\textwidth]{zetun_accueil.png}
  \caption{Interface premium de la plateforme Zetun}
  \label{fig:zetun_accueil}
\end{figure}

\begin{figure}[H]
  \centering
  \includegraphics[width=\textwidth]{zetun_produits.png}
  \caption{Présentation des produits sur la plateforme Zetun}
  \label{fig:zetun_produits}
\end{figure}

% ---- Plateforme 3 : Oleavanti ----
\subsubsection*{Oleavanti}
% [Description]

\begin{figure}[H]
  \centering
  \includegraphics[width=\textwidth]{oleavanti_accueil.png}
  \caption{Interface d'accueil de la plateforme Oleavanti}
  \label{fig:oleavanti_accueil}
\end{figure}

\begin{figure}[H]
  \centering
  \includegraphics[width=\textwidth]{oleavanti_produits.png}
  \caption{Présentation des produits sur la plateforme Oleavanti}
  \label{fig:oleavanti_produits}
\end{figure}

% ---- Plateforme 4 : Castillo de Canena (internationale) ----
\subsubsection*{Castillo de Canena (Espagne)}
% [Description]

\begin{figure}[H]
  \centering
  \includegraphics[width=\textwidth]{castillo_accueil.png}
  \caption{Interface d'accueil de la plateforme Castillo de Canena}
  \label{fig:castillo_accueil}
\end{figure}

\begin{figure}[H]
  \centering
  \includegraphics[width=\textwidth]{castillo_produits.png}
  \caption{Présentation des produits sur la plateforme Castillo de Canena}
  \label{fig:castillo_produits}
\end{figure}

% ---- Plateforme 5 : Frantoio Muraglia (internationale) ----
\subsubsection*{Frantoio Muraglia (Italie)}
% [Description]

\begin{figure}[H]
  \centering
  \includegraphics[width=\textwidth]{muraglia_accueil.png}
  \caption{Interface d'accueil de la plateforme Frantoio Muraglia}
  \label{fig:muraglia_accueil}
\end{figure}

\begin{figure}[H]
  \centering
  \includegraphics[width=\textwidth]{muraglia_produits.png}
  \caption{Présentation des produits sur la plateforme Frantoio Muraglia}
  \label{fig:muraglia_produits}
\end{figure}

\subsection{Critique de l'existant}

\begin{table}[H]
\centering
\caption{Tableau comparatif des plateformes étudiées (national et international)}
\label{tab:existant}
\begin{tabularx}{\textwidth}{>{\bfseries}p{3.5cm} X X}
\toprule
\rowcolor{beige}
\textbf{Plateforme} & \textbf{Points forts} & \textbf{Points faibles} \\
\midrule
Carthage Crown
& Design élégant, valorisation du terroir tunisien, navigation fluide, identité visuelle premium
& Absence de visualisation 3D, interaction utilisateur limitée, expérience principalement statique \\
\midrule
ZETUN
& Interface moderne, animations fluides, branding cohérent, esthétique haut de gamme
& Absence de configurateur interactif, aucune personnalisation produit, immersion limitée \\
\midrule
Oleavanti
& Storytelling visuel fort, univers graphique luxueux, mise en valeur émotionnelle du produit
& Présentation basée sur des images fixes, absence d'exploration immersive, manque d'interactivité avancée \\
\midrule
Castillo de Canena \newline \textcolor{gris}{\small Espagne}
& Design luxueux, storytelling familial fort, e-commerce international complet, packaging premium mis en valeur
& Absence de configurateur interactif, présentation produit statique, pas de visualisation 3D \\
\midrule
Frantoio Muraglia \newline \textcolor{gris}{\small Italie}
& Identité visuelle artisanale raffinée, univers de marque cohérent, catalogue produits détaillé
& Aucune personnalisation produit en ligne, expérience purement informative, pas d'interactivité avancée \\
\bottomrule
\end{tabularx}
\end{table}

\noindent Cette analyse élargie à des références internationales confirme qu'aucun acteur du secteur oléicole premium, qu'il soit tunisien, espagnol ou italien, ne propose à ce jour de configurateur interactif à visualisation 3D — ce qui constitue la principale valeur ajoutée et l'axe de différenciation central de la plateforme FENDRI.

% -----------------------------------------------------------
\section{Solution proposée}
% -----------------------------------------------------------
% [Copier le texte depuis Canva]

% -----------------------------------------------------------
\section{Méthodologie de développement}
% -----------------------------------------------------------

\subsection{Approche agile de développement}

\begin{figure}[H]
  \centering
  \includegraphics[width=0.7\textwidth]{cycle_agile.png}
  \caption{Cycle de développement agile}
  \label{fig:cycle_agile}
\end{figure}

\subsection{Étude comparative des méthodes}

\begin{table}[H]
\centering
\caption{Tableau comparatif des méthodes de développement}
\label{tab:methodes}
\begin{tabularx}{\textwidth}{>{\bfseries}p{2cm} X X X}
\toprule
\rowcolor{beige}
\textbf{Méthode} & \textbf{Description} & \textbf{Avantages} & \textbf{Limites} \\
\midrule
RUP    & Méthode structurée basée sur un développement itératif & Bonne organisation et documentation & Complexe à mettre en œuvre \\
\midrule
2TUP   & Méthode séparant les aspects fonctionnels et techniques & Bonne séparation des besoins et de l'architecture & Nécessite une planification importante \\
\midrule
Scrum  & Méthode agile basée sur des cycles courts appelés sprints & Flexible, collaborative et adaptée aux changements & Dépend fortement de l'organisation de l'équipe \\
\midrule
Kanban & Méthode agile basée sur la gestion visuelle des tâches & Facilite le suivi du travail et améliore la productivité & Peut manquer de structure pour les grands projets \\
\midrule
XP     & Méthode agile axée sur la qualité du code & Améliore la qualité logicielle et favorise les retours rapides & Demande une forte implication de l'équipe et du client \\
\midrule
SAFe   & Cadre agile pour les projets à grande échelle & Permet de coordonner plusieurs équipes agiles efficacement & Mise en œuvre complexe et coûteuse \\
\bottomrule
\end{tabularx}
\end{table}

\subsection{Présentation de la méthode Scrum}
% [Copier le texte depuis Canva]

\begin{figure}[H]
  \centering
  \includegraphics[width=0.8\textwidth]{cycle_scrum.png}
  \caption{Cycle de fonctionnement de la méthodologie Scrum}
  \label{fig:cycle_scrum}
\end{figure}

\begin{figure}[H]
  \centering
  \includegraphics[width=0.7\textwidth]{roles_scrum.png}
  \caption{Les rôles Scrum du projet FENDRI}
  \label{fig:roles_scrum}
\end{figure}

\subsection{Artefacts de Scrum}
% [Copier le texte depuis Canva]

\subsection{Avantages de la méthode Scrum}
% [Copier le texte depuis Canva]

\subsection{Analyse des risques du projet}
% [Tableau des risques — copier depuis Canva]

\subsection{Estimation des ressources}
% [Tableaux ressources humaines et techniques]

\subsection{WBS — Work Breakdown Structure}

\begin{figure}[H]
  \centering
  \includegraphics[width=\textwidth]{wbs.png}
  \caption{WBS — Work Breakdown Structure}
  \label{fig:wbs}
\end{figure}

\subsection{Jalons du projet}
% [Copier le texte depuis Canva]

\subsection{Backlog de Sprint}
% [Tableaux backlog]

\subsection{Planification des sprints}
% [Tableaux sprints 1 à 4]

\subsection{Diagramme de Gantt}

\begin{figure}[H]
  \centering
  \includegraphics[width=\textwidth]{gantt.png}
  \caption{Diagramme de Gantt}
  \label{fig:gantt}
\end{figure}

% -----------------------------------------------------------
\section{Pertinence du sujet}
% -----------------------------------------------------------
% [Copier le texte depuis Canva]

% -----------------------------------------------------------
\section{Différenciation concurrentielle et innovation digitale}
% -----------------------------------------------------------
% [Copier le texte depuis Canva]

\section*{Conclusion}
\addcontentsline{toc}{section}{Conclusion}
Dans ce chapitre, j'ai présenté le cadre général du projet ainsi que le contexte de développement de la plateforme FENDRI. J'ai également exposé les objectifs du projet, la problématique identifiée ainsi qu'une étude de l'existant permettant de mieux comprendre les besoins et les limites des solutions actuelles. Le chapitre suivant sera consacré à la conception 3D des produits et à l'identité visuelle de la marque FENDRI.
